% Draft (venue to be decided): annealing simulators vs D-Wave Advantage. Anonymous.
\documentclass[11pt]{article}
\usepackage[T1]{fontenc}
\usepackage[utf8]{inputenc}
\usepackage[margin=2.5cm]{geometry}
\usepackage{amsmath,amssymb}
\usepackage{booktabs}
\usepackage{graphicx}
\usepackage{url}
\usepackage[hidelinks,pdfauthor={},pdftitle={}]{hyperref}
\newcommand{\tofill}[1]{\textbf{[#1]}}

\title{Can Simulators Stand In for a Quantum Annealer? Classical, Quantum
Monte Carlo and Exact Simulation against D-Wave Advantage in Neighborhood
Search}
\author{Anonymous}
\date{}

\begin{document}
\maketitle

\begin{abstract}
\tofill{Draft: the abstract is written after the results.}
\end{abstract}

\section{Introduction}

Quantum annealers solve quadratic unconstrained binary optimization (QUBO)
problems by slowly turning off a transverse field~\cite{kadowaki1998,albash2018aqc},
and commercial devices with thousands of qubits are available~\cite{johnson2011manufactured}.
Access to them is metered, so researchers often replace the device with a
simulator during development: classical simulated annealing, simulated quantum
annealing by path-integral Monte Carlo, or, for very small systems, the exact
dynamics. Earlier studies compared such simulators with D-Wave devices on
random spin glasses and asked whether the device behaves like a quantum or a
classical annealer~\cite{boixo2014evidence,ronnow2014speedup}. Practitioners ask
a narrower question: if a simulator replaces the device inside an optimization
algorithm, does the algorithm reach the same result?

We answer it for a concrete workload: neighborhood search for the permutation
flow shop problem, where a quantum annealer selects the move of iterated local
search by solving a windowed
QUBO~\cite{wojewodzki2026windowed,wojewodzki2026bpasts}. A large set of runs on
D-Wave Advantage already exists, so the comparison needs no new hardware time.
We run the same QUBOs on several simulators under three scenarios of
increasing realism, compare the makespans they reach with those of the device
and report the cost of each simulator. \tofill{One sentence with the main
finding once the results are in.}

\section{Background}

\paragraph{Workload.}
Iterated local search~\cite{lourenco2003} for the permutation flow shop
problem on the Taillard instances~\cite{taillard1993} moves to the best
permutation in an exponential neighborhood: Adjacent, Fibonacci~\cite{bozejko2026new},
Dynasearch or Motzkin~\cite{wojewodzki2026soco}. The move is the solution of a
QUBO~\cite{lucas2014ising,glover2019} whose variables select swaps; for large
permutations the search decomposes the problem into overlapping windows and
solves each on the annealer~\cite{wojewodzki2026windowed}. Earlier work used
D-Wave devices for single-machine scheduling as well~\cite{bozejko2022distributed}.

\paragraph{The device.}
D-Wave Advantage (system 4.1, Pegasus topology) programs the Ising
coefficients $h_i$ and $J_{ij}$ of the problem on superconducting flux qubits.
A logical variable whose couplings do not fit the hardware graph becomes a
chain of physical qubits held together by a strong ferromagnetic coupling, a
minor embedding~\cite{choi2008minor}; chains can break, and the solution then
needs repair. The programmed coefficients carry integrated control errors,
Gaussian deviations of $h$ and $J$ with a systematic and a random
part~\cite{dwaveice}, and the device anneals at a finite temperature. Recent
experiments show coherent quantum dynamics in programmable Ising
chains~\cite{king2022coherent}.

\section{Simulators}

\paragraph{Exact ground state.}
For windows of up to about twenty variables, enumeration gives the optimum of
every QUBO. It shows what a perfect annealer would return and is the
reference for the others.

\paragraph{Simulated annealing.}
Simulated annealing~\cite{kirkpatrick1983sa} samples the Boltzmann
distribution at a decreasing temperature with Metropolis updates. D-Wave's
own classical samplers~\cite{dwavesamplers} implement it, together with tabu
search, for the same problem format as the device, which makes it the
natural classical stand-in.

\paragraph{Simulated quantum annealing.}
Simulated quantum annealing samples the transverse-field Ising model by
path-integral Monte Carlo: the quantum system maps to a classical one with an
extra imaginary-time dimension of Trotter slices, and the transverse field is
lowered during the run~\cite{santoro2002qa,martonak2002pimc}. It reproduces the
equilibrium statistics of quantum annealing, not its real-time dynamics, and it
was the reference against which early D-Wave devices were
judged~\cite{boixo2014evidence}. OpenJij~\cite{openjij} provides both simulated
annealing and simulated quantum annealing with configurable schedules.

\paragraph{Exact dynamics.}
QuTiP~\cite{johansson2013qutip,lambert2024qutip5} integrates the
Schr\"{o}dinger or master equation of the annealing Hamiltonian along a given
schedule, which includes decoherence explicitly. The cost grows exponentially
with the number of qubits, so we use it on a few small windows only.

\paragraph{Quantum-inspired solvers.}
Simulated bifurcation~\cite{goto2019sb} and the digital
annealer~\cite{aramon2019da} solve QUBOs with dynamics inspired by physical
systems. They do not simulate a quantum annealer, and part of them is
commercial, but they are common replacements for it; we consider them as
optional baselines.

\begin{table}[t]
\caption{Annealing simulators. The last two columns come from our measurements.}
\label{tab:sims}
\centering
\small
\begin{tabular}{lllcc}
\toprule
Simulator & What it computes & Availability & max variables & s / window \\
\midrule
Enumeration & exact ground state & own code & \tofill{?} & \tofill{?} \\
dwave-samplers (SA, tabu) & classical annealing & open source & \tofill{?} & \tofill{?} \\
OpenJij (SA, SQA) & classical and quantum Monte Carlo annealing & open source & \tofill{?} & \tofill{?} \\
QuTiP & exact (open-system) dynamics & open source & \tofill{?} & \tofill{?} \\
Simulated bifurcation, digital annealer & quantum-inspired heuristics & partly commercial & \tofill{?} & \tofill{?} \\
\bottomrule
\end{tabular}
\end{table}

\section{Methodology}

\paragraph{Scenarios.}
In A1 the simulator returns the exact ground state of each window. In A2 it
runs annealing dynamics without control errors: simulated annealing and
simulated quantum annealing with a schedule matched to the device (anneal time,
effective temperature), and exact dynamics on small windows. A3 adds the
device's control errors as Gaussian perturbations of $h$ and $J$ and solves the
problem through the same minor embedding as the device, so chains can break.

\paragraph{Comparison with the device.}
The device runs did not store full sample sets, so the comparison works at the
level of whole runs: each simulator replaces the device inside iterated local
search on the same instances and windows, with the same number of annealer
calls the device made, and we compare the relative deviation of the makespan
from the Taillard lower bound with a paired sign test and Kendall's $\tau$
over instances. If annealer time allows, a small run that stores the sample
sets of selected windows gives a distribution-level comparison \tofill{?}.

\section{Experimental Setup}

The study runs in three stages: every simulator on one window; all of them on
the instances with $n = 20$; the most faithful ones on the full set
($n = 20$ to $500$), which simulated annealing and simulated quantum annealing
can afford.

\begin{table}[t]
\caption{D-Wave data used as reference (D-Wave Advantage system 4.1, Pegasus).}
\label{tab:hw}
\centering
\small
\begin{tabular}{ll}
\toprule
Item & Value \\
\midrule
Runs (ILS / SA) & 5650 (3250 / 2400), May--August 2026 \\
Instances & $n = 20$: 5200 runs; $n = 50$: 200; $n = 100$: 100; $n = 200$: 100; $n = 500$: 50 \\
Neighborhoods & Adjacent, Fibonacci, Dynasearch, Motzkin; basic and enhanced variants~\cite{wojewodzki2026windowed} \\
Annealer calls & 24\,317 (11 May -- 8 August 2026) \\
QPU access time & 753 s in total \\
Problem size per call & median 16 variables, maximum 220 \\
Stored per run & makespan history, best permutation, call statistics (no sample sets) \\
\bottomrule
\end{tabular}
\end{table}

\section{Results}

\tofill{Draft: written after the three stages.}

\begin{figure}[t]
\centering
\fbox{\parbox[c][3.5cm][c]{0.9\textwidth}{\centering \tofill{Figure: relative
deviation of the makespan, simulator versus D-Wave, one point per instance, for
A1, A2 (SA, SQA) and A3.}}}
\caption{Run-level agreement with the device.}
\label{fig:prd}
\end{figure}

\begin{table}[t]
\caption{Median relative deviation [\%] at the same number of annealer calls as the device ($n = 20$).}
\label{tab:prd}
\centering
\small
\begin{tabular}{lccccc}
\toprule
Neighborhood & D-Wave & A1 exact & A2 SA & A2 SQA & A3 SQA + errors \\
\midrule
Adjacent & \tofill{?} & \tofill{?} & \tofill{?} & \tofill{?} & \tofill{?} \\
Fibonacci & \tofill{?} & \tofill{?} & \tofill{?} & \tofill{?} & \tofill{?} \\
Dynasearch & \tofill{?} & \tofill{?} & \tofill{?} & \tofill{?} & \tofill{?} \\
Motzkin & \tofill{?} & \tofill{?} & \tofill{?} & \tofill{?} & \tofill{?} \\
\bottomrule
\end{tabular}
\end{table}

\section{Discussion}
\tofill{Draft: which simulator can stand in for the device and for what;
limitations: no sample sets from the device, error model from documentation,
one device generation.}

\section{Conclusions}
\tofill{Draft.}

\bibliographystyle{plain}
\bibliography{references}
\end{document}
